\documentclass{article}
\usepackage[utf8]{inputenc}
\usepackage[T1]{fontenc}
\usepackage{geometry}
\usepackage{graphicx}
\usepackage{float}
\usepackage{hyperref}
\usepackage[english,indonesian]{babel}

\title{Optimasi Pemrosesan Template LaTeX Menggunakan Pendekatan Berbasis AI}
\author{Academic Support Agent}
\date{\today}

\begin{document}
\maketitle

\begin{abstract}

\end{abstract}

\chapter{PENDAHULUAN}


\chapter{TINJAUAN PUSTAKA}


% Original content (commented out):
% \documentclass[a4paper]{article}

\usepackage[english]{babel}
\usepackage[utf8x]{inputenc}
\usepackage{amsmath}
\usepackage{graphicx}
\usepackage[colorinlistoftodos]{todonotes}

\usepackage{siunitx}
\usepackage{booktabs}
\usepackage{subcaption}
\usepackage[section]{placeins}

\usepackage{tikz}
\usetikzlibrary{patterns}

\usepackage[margin=1in]{geometry}

\setlength{\parindent}{0cm}

\title{CFD Template}

\author{Reginald Pierce}

\date{\today}

\begin{document}
\maketitle
\vspace{-1.2in}
\includegraphics[width=2.2cm]{author.jpg}
\setlength{\parskip}{0.2cm}

\newpage
\section{Problem Definition}

\begin{figure}[h]
	\centering
    \begin{tikzpicture}
    	\draw[thick] (0,0) -- (1,0) -- (1,1) -- (0,1) -- cycle;
    \end{tikzpicture}
    \caption{Problem Schematic}
    \label{fig:problem}
\end{figure}

\section{Mesh}

\section{Numerical Procedures}

\section{Results}

\section{Conclusion}



\end{document}

\end{document}